%==========================================================================================
\section{Follow-up (21 October): fair-horizon training (2 $\to$ 8)}\label{sec:fairh8}

The requirements for this round were:
\begin{itemize}
\item the model is trained \textbf{only for forecasting}, with no data-assimilation objective;
\item it is trained \textbf{end-to-end}, with no frozen or staged parts;
\item its training horizon is \textbf{comparable to the baselines}: 2 conditioning frames and 8
  predicted frames, as in U-Net\_ut's 8-frame unroll and the original KAE's 2 $\to$ 8;
\item its latent size is accurate but affordable.
\end{itemize}

\textbf{Setting} (\texttt{configs/ctkae/fair\_h8.yaml}):
\begin{itemize}
\item every one of the 8 predicted frames is supervised, at stride 2 only ($\Delta t = 0.1$, as
  ACDM);
\item ACDM's 33 training simulations, all frames, batch 16, 20k steps, final checkpoint, no
  validation split;
\item purely linear latent dynamics $dz/dt = A(Ma)\,z$, with no $\mu$ (\S\ref{sec:mu}).
\end{itemize}
Compared with the 60-step model of \S\ref{sec:exp}, this round also changes the training data
(33 vs 30 simulations), the number of steps (20k vs 12k) and the supervision (every frame vs
random frames). The gain below therefore cannot be credited to the short horizon alone.

\subsection{Forecasting}

\begin{center}\small
\begin{tabular}{lcccccc}
\toprule
Test, MSE $\times10^{-3}$ & interp & extrap & long & lead 1 & energy at 58 & corr.\ at 58\\
\midrule
\textbf{H8, latent 256} (3 seeds) & 3.49 $\pm$ 0.15 & 1.98 $\pm$ 0.31 & 16.2 $\pm$ 0.6 & 0.26 & \textbf{0.94} & 0.69\\
H8, latent 128 (seed 0)           & 3.84 & 1.92 & 15.9 & 0.35 & 0.89 & --\\
2 $\to$ 60 (3 seeds, \S\ref{sec:exp}) & 4.07 $\pm$ 0.29 & 2.18 $\pm$ 0.14 & 16.5 $\pm$ 2.0 & 0.66 & 0.77 & 0.62\\
ACDM (official)                   & 2.81--3.02 & 1.99--2.22 & 22.5--22.8 & 0.03 & 1.10 & 0.73\\
U-Net (official, best variant)    & 1.53 & 1.42 & 22.4 & 0.02 & 1.08 & 0.85\\
\bottomrule
\end{tabular}
\end{center}
{\footnotesize Per seed, latent 256: interp 3.43 / 3.38 / 3.66, extrap 1.67 / 1.97 / 2.29,
long 16.8 / 15.6 / 16.2.}

\textbf{Result.}
\begin{itemize}
\item \textbf{This is the best CT-KAE so far.} Compared with the 60-step model:
  \begin{itemize}
  \item interpolation error is 14\% lower and extrapolation error 9\% lower;
  \item the first-step error is 2.5$\times$ lower (0.26 vs 0.66);
  \item it keeps 94\% of the oscillation energy at step 58 instead of 77\%, so the
    mid-rollout damping is mostly gone.
  \end{itemize}
\item \textbf{Against ACDM:}
  \begin{itemize}
  \item it is still behind on interpolation (3.49 vs about 2.9);
  \item it is level on extrapolation (1.98 vs 1.99--2.22; seed spread 1.67--2.29);
  \item it has lower 240-step error.
  \end{itemize}
  The U-Nets remain better on both 60-step cases.
\item \textbf{Latent size.} 256 beats 128 on interpolation, the first-step error, the energy at
  step 58 and every DA number. On extrapolation and the 240-step case the two are within the
  seed spread. 256 costs 24.7\,M parameters against 16.3\,M, at the same training time: about
  1 hour and 12.3\,GB on one GPU when not shared. 256 is the recommended size.
\item \textbf{Zero-shot still holds.} Predictions at unseen half-steps are as accurate as those
  at whole steps (3.450 vs 3.434), and the response to Mach number is smooth.
\end{itemize}

\subsection{Data assimilation}

\textbf{Setup.}
\begin{itemize}
\item \textbf{Same pipeline as \S\ref{sec:da}:} 4D-Var code, KAE slot, and learning rate tuned
  on \texttt{gt\_extrap}.
\item \textbf{Wider learning-rate grid.} The earlier grid (0.03--1.0) ended at its top value
  for the new models, while the error was still falling. So I widened it to 0.3--10. The
  optimum is now inside the grid at 3.0, for both latent sizes and all three seeds (tuning
  error 0.034 at 3.0, against 0.043 at 1.0 and 0.053 at 10).
\end{itemize}

Mean relative error over each sweep's points (8 problems per point; $\pm$ is the standard
deviation over seeds):

\begin{center}\footnotesize
\begin{tabular}{lccccc}
\toprule
Sweep & Original & H8, 256 & H8, 256 & H8, 128 & Curriculum\\
 & KAE & lr 3, 3 seeds & lr 1, seed 0 & lr 1, seed 0 & lr 1\\
\midrule
\multicolumn{6}{l}{\emph{df = 9}}\\
window length $\delta_l$ & \textbf{0.085} & 0.096 $\pm$ 0.002 & 0.100 & 0.109 & 0.094\\
noise                    & \textbf{0.078} & 0.094 $\pm$ 0.002 & 0.100 & 0.110 & 0.094\\
sparsity                 & \textbf{0.086} & 0.098 $\pm$ 0.002 & 0.104 & 0.115 & 0.100\\
misspecified noise       & \textbf{0.078} & 0.093 $\pm$ 0.002 & 0.100 & 0.110 & 0.099\\
number of observations   & \textbf{0.079} & 0.095 $\pm$ 0.004 & 0.100 & 0.111 & 0.095\\
joint sparsity           & \textbf{0.078}$^\ast$ & 0.097 $\pm$ 0.005 & 0.101 & 0.112 & 0.100\\
\midrule
\multicolumn{6}{l}{\emph{df = 1 (seed 0 only for H8)}}\\
first-observation lead $\delta_f$ & \textbf{0.070} & 0.076 & 0.083 & 0.093 & 0.074\\
single observation       & 0.089 & \textbf{0.080} & 0.088 & 0.093 & 0.091\\
window length $\delta_l$ & \textbf{0.027} & 0.040 & 0.047 & 0.055 & 0.027\\
noise                    & \textbf{0.028} & 0.041 & 0.049 & 0.058 & 0.028\\
sparsity                 & \textbf{0.031} & 0.047 & 0.056 & 0.069 & 0.037\\
misspecified noise       & \textbf{0.028} & 0.041 & 0.049 & 0.058 & 0.029\\
number of observations   & \textbf{0.026} & 0.042 & 0.048 & 0.059 & 0.027\\
joint sparsity           & \textbf{0.027} & 0.043 & 0.050 & 0.062 & 0.028\\
\midrule
\multicolumn{6}{l}{\emph{Representation check (df = 9 default problem)}}\\
decoder floor            & 0.041 & \textbf{0.026} & 0.027 & 0.031 & 0.032\\
4D-Var started from the true state & \textbf{0.074} & 0.088 & 0.086 & 0.096 & 0.096\\
\bottomrule
\end{tabular}
\end{center}
{\footnotesize $^\ast$ 3 of 5 points. The curriculum model was not re-tuned on the wider grid;
its learning rate is also at the old grid top, so its numbers may also improve.}

\textbf{Result.}
\begin{itemize}
\item \textbf{The original KAE is still the best data-assimilation model.}
  \begin{itemize}
  \item df = 9: the fair-horizon model is 13--24\% behind it (0.093--0.098 vs
    0.078--0.086). The gap is consistent across the three seeds.
  \item df = 1: it is about 1.5$\times$ the original's error on the window sweeps, but slightly
    better on the single-observation sweep.
  \end{itemize}
\item \textbf{The fair-horizon model is level with the curriculum model at df = 9.} It is
  worse at df = 1.
\item \textbf{Tuning mattered.} The wider learning-rate grid alone improved the new model
  by 7--10\% at df = 9 and by 9--16\% at df = 1 (seed 0).
\item \textbf{The gap is not representation.} The new decoder reproduces the true state
  better than the original's (floor 0.026 vs 0.041). The gap opens when 4D-Var fits
  observations 9--25 steps ahead (0.088 vs 0.074 from the true state).
\item \textbf{Forecast skill does not explain it either.} Over leads 9--25 the new model's
  forecast error is 3.4$\times$ lower than the original's on interpolation (1.14 vs 3.83) and
  2.1$\times$ lower on extrapolation (0.66 vs 1.36). Two causes remain untested:
  \begin{itemize}
  \item the original DA checkpoint was selected on the regime DA is tested on;
  \item the latent coordinates of the two models are conditioned differently for the inverse
    problem.
  \end{itemize}
  A DA-style training objective, which might close the gap, is excluded by design: the model
  is trained only for forecasting.
\end{itemize}
